% Chapter 4, Figure 5(b): percent error in y_b, B14 engine (scan: figures/ch4/fig05b.png).
% Curves: digitized from the 1973 art by fig05b.py (the book gives neither the B14 thrust curve nor its
% propellant mass, so its interval solution (83)-(87) cannot be rerun); template fig06a.tex. As printed, the
% two Fehskens-Malewicki curves are one line, and the Caporaso-Bengen k_min curve joins it near 0.05.
\documentclass[11pt]{standalone}
\usepackage{tamrfig}
\begin{document}
\begin{tikzpicture}
  \begin{axis}[tamr,
      width=4.0in, height=2.5in,
      xmin=0.02, xmax=0.14, xtick={0.02,0.04,0.06,0.08,0.10,0.12,0.14},
      minor xtick={0.03,0.05,0.07,0.09,0.11,0.13},
      xticklabel style={/pgf/number format/.cd, fixed, fixed zerofill, precision=2},
      ymin=-15, ymax=15, ytick distance=5,
      xlabel={$m_o$ (kg)}, ylabel={Percent error in $y_b$},
      legend pos=south east, legend style={nodes={inner sep=1pt}}]
    \draw[muted, line width=0.4pt] (axis cs:0.02,0) -- (axis cs:0.14,0);
    \addplot[series1] table[x=m0, y=fm_kmax, col sep=comma] {figures/v2/ch4/fig05b.csv};
    \addlegendentry{Fehskens-Malewicki}
    \addplot[series2] table[x=m0, y=cb_kmax, col sep=comma] {figures/v2/ch4/fig05b.csv};
    \addlegendentry{Caporaso-Bengen}
    \addplot[series1] table[x=m0, y=fm_kmin, col sep=comma] {figures/v2/ch4/fig05b.csv};
    \addplot[series2] table[x=m0, y=cb_kmin, col sep=comma] {figures/v2/ch4/fig05b.csv};
    % k labels with a leader to each method's curve for that k (points on the fig05b.csv curves)
    \node (kmax) at (axis cs:0.0455,9.0) {\footnotesize $k_{\max}$};
    \draw[leader] (kmax.west) -- (axis cs:0.0220,8.247);
    \draw[leader] (kmax.south west) -- (axis cs:0.0350,-0.287);
    \node (kmin) at (axis cs:0.0265,2.9) {\footnotesize $k_{\min}$};
    \draw[leader] (kmin.north west) -- (axis cs:0.0235,8.064);
    \draw[leader] (kmin.north east) -- (axis cs:0.0310,6.167);
  \end{axis}
\end{tikzpicture}
\end{document}
